\subsection*{Elementos que componen la infraestructura}
\addcontentsline{toc}{subsection}{Elementos que componen la infraestructura}

Una infraestructura en la nube no se entiende únicamente como un conjunto de recursos aislados, sino como una arquitectura coherente donde cada componente cumple una función específica dentro del sistema y se integra con los demás para formar un entorno operativo completo. El propósito de esta sección es describir los distintos elementos que conforman la infraestructura seleccionada para el proyecto \textbf{Muncher}, analizando su papel en la arquitectura general y justificando las decisiones técnicas que motivan su adopción.

\subsubsection*{Front-end}
\addcontentsline{toc}{subsubsection}{Front-end}

Para el componente de \textbf{front‑end}, al tratarse de una aplicación de página única (\textit{Single Page Application}, SPA), fue necesario decidir cómo publicar los archivos resultantes de su compilación. \textbf{AWS} ofrece varias vías para ello: el alojamiento de sitios estáticos sobre \textbf{Amazon S3}, la combinación de dicho almacenamiento con la red de distribución \textbf{Amazon CloudFront}, o servicios gestionados de mayor nivel como \textit{AWS Amplify Hosting}. Las tres alternativas resultan válidas para servir contenido estático, pero difieren en el grado de control sobre la configuración, en el coste y en la capacidad de distribución geográfica.

Amazon S3 permite alojar un sitio web estático a partir de los archivos depositados en un \textit{bucket}, entendiendo por estático aquel cuyas páginas incluyen contenido fijo y, opcionalmente, scripts ejecutados en el cliente.\footnote{Documentación oficial: \url{https://docs.aws.amazon.com/AmazonS3/latest/userguide/WebsiteHosting.html}} Sobre este almacenamiento se sitúa \textbf{Amazon CloudFront}, que distribuye el contenido a través de una red mundial de centros de datos denominados \textit{edge locations}: cada petición se encamina a la ubicación que ofrece la menor latencia y, si el archivo no está en ella, se recupera del origen definido —en este caso el \textit{bucket} de S3— y se almacena en caché para las siguientes.\footnote{Documentación oficial: \url{https://docs.aws.amazon.com/AmazonCloudFront/latest/DeveloperGuide/Introduction.html}} CloudFront asigna además un dominio propio a cada distribución, admite el uso de un nombre de dominio alternativo y permite controlar el tiempo de permanencia de los archivos en caché. La transferencia de datos entre el origen y CloudFront no se factura cuando el origen es un servicio de AWS como S3, de modo que el coste se limita a la salida de datos hacia el usuario y al número de peticiones atendidas.

\begin{figure}[htbp]
	\centering
	\includegraphics[width=0.85\textwidth]{front-end-infrastructure.pdf}
	\caption{Infraestructura del front-end.}
	\label{fig:front-end-infrastructure}
\end{figure}

Se optó por la combinación de \textbf{S3 y CloudFront} porque el front‑end de \textbf{Muncher} no requiere ejecución de código en el servidor: basta con almacenar y servir los archivos del cliente compilado (HTML, CSS, JavaScript y recursos estáticos), lo que elimina la necesidad de instancias o procesos persistentes y reduce el coste de la infraestructura. Se descartó \textit{AWS Amplify Hosting} pese a ser la opción que la propia documentación de S3 recomienda para este caso, ya que su carácter de servicio totalmente gestionado abstrae buena parte de la configuración subyacente y entra en conflicto con el criterio de alta capacidad de configuración establecido al inicio de esta sección. Frente a ello, S3 y CloudFront ofrecen el mismo resultado —entrega global mediante CDN de una aplicación cliente— conservando el control explícito sobre el almacenamiento, la política de caché y la distribución.

\subsubsection*{API}
\addcontentsline{toc}{subsubsection}{API}

Para el componente de \textbf{API} la decisión es de naturaleza distinta a la del front-end: no se trata de una aplicación estática, sino que requiere cómputo para atender cada petición, de modo que la elección recae sobre el modelo de ejecución. \textbf{AWS} ofrece dos familias de opciones para ello. La primera son los contenedores gestionados —\textbf{Amazon ECS} sobre \textit{Fargate}, \textbf{AWS App Runner} o los servicios de contenedores de \textbf{Amazon Lightsail}—, en los que la aplicación se empaqueta en una imagen y se mantiene un proceso en ejecución de forma permanente. La segunda son las funciones de \textbf{AWS Lambda}, donde el proveedor ejecuta el código únicamente durante el tiempo que tarda en atender cada invocación. A esta elección se añade una segunda, independiente de la anterior: cómo exponer la API a Internet, para lo que caben \textbf{Amazon API Gateway}, un balanceador de carga de aplicación o la propia URL de función situada detrás de la distribución de CloudFront descrita en el apartado anterior.

El precio es lo que separa con mayor claridad a ambas familias. Un contenedor en ejecución continua se cobra por horas con independencia de que reciba tráfico: la tarea más pequeña que admite \textit{Fargate} —0,25 vCPU y 0,5 GB de memoria— supone, a los precios publicados para la región de Norte de Virginia, un gasto próximo a los 9 USD mensuales si permanece activa todo el mes,\footnote{Precios oficiales: \url{https://aws.amazon.com/ecs/pricing/}} y el plan más reducido de los servicios de contenedores de \textbf{Lightsail} asciende a 7 USD mensuales fijos.\footnote{Precios oficiales: \url{https://aws.amazon.com/lightsail/pricing/}} A ello se suma que \textbf{App Runner} dejó de admitir nuevos clientes el 30 de abril de 2026, por lo que ha dejado de ser una alternativa disponible.\footnote{Aviso oficial: \url{https://aws.amazon.com/apprunner/}} \textbf{Lambda}, por el contrario, no factura mientras no se invoca, y su capa gratuita —un millón de peticiones y 400.000 GB-segundo de cómputo al mes— es permanente y no expira a los doce meses como ocurre con otras.\footnote{Precios oficiales: \url{https://aws.amazon.com/lambda/pricing/}} Con el tráfico previsto para este sistema, el coste de la API es por tanto nulo, frente a un gasto fijo que por sí solo agotaría el margen que fija el criterio de coste reducido establecido al inicio de esta sección.

Se optó en consecuencia por \textbf{AWS Lambda}, y la segunda decisión se resolvió reutilizando la distribución de \textbf{CloudFront} que ya sirve el front-end. A ella se añadió un comportamiento de caché para las rutas \texttt{/api/*} cuyo origen es la URL de función, de manera que la API se publica en el mismo dominio que la interfaz y bajo esa ruta. Se descartó \textbf{API Gateway} porque su cometido —enrutar las peticiones y publicar la API— lo cumple ya la distribución existente, y porque su capa gratuita de un millón de llamadas cubre únicamente los doce primeros meses desde la creación de la cuenta,\footnote{Precios oficiales: \url{https://aws.amazon.com/api-gateway/pricing/}} de forma que incorporarlo equivaldría a programar un gasto futuro sin obtener ninguna capacidad que el sistema no tenga. El balanceador de carga se descartó por la misma razón que los contenedores: se factura por horas de existencia y no por uso.

\begin{figure}[htbp]
	\centering
	\includegraphics[width=0.85\textwidth]{api-infrastructure.pdf}
	\caption{Infraestructura de la API.}
	\label{fig:api-infrastructure}
\end{figure}

El modelo elegido impone dos limitaciones que conviene hacer explícitas. La primera es el arranque en frío: cuando una invocación no encuentra un entorno de ejecución ya disponible, al tiempo de respuesta se añade la inicialización del intérprete y la importación de las dependencias de la aplicación. La segunda afecta a las peticiones con cuerpo, ya que Lambda no admite cargas sin firmar cuando el origen está protegido de esta manera: los métodos POST y PUT exigen que el cliente calcule el SHA-256 del cuerpo y lo envíe en la cabecera \texttt{x-amz-content-sha256}.

\subsubsection*{Monitorización y observabilidad}
\addcontentsline{toc}{subsubsection}{Monitorización y observabilidad}

Para el componente de \textbf{monitorización} la decisión no versa sobre dónde se ejecuta el sistema, como en los dos apartados anteriores, sino sobre dónde se recogen, se almacenan y se consultan los datos que describen su funcionamiento. Se distinguen tres clases de datos, y los tres son necesarios para satisfacer los requerimientos de esta categoría: los \textit{logs}, que dejan escritos los sucesos que la aplicación considera relevantes; las trazas, que siguen una misma petición a través de los componentes que intervienen en atenderla y miden el tiempo consumido en cada uno; y las métricas, que agregan el comportamiento del sistema en series temporales.

\textbf{AWS} cubre este cometido con \textbf{Amazon CloudWatch} —que reúne el almacenamiento y la consulta de \textit{logs}, las métricas, los paneles y las alarmas— y con \textbf{AWS X-Ray}, que recoge las trazas. Fuera del proveedor existe un conjunto de plataformas de observabilidad cuyas capas gratuitas son considerablemente más amplias: \textbf{Grafana Cloud} ofrece 50 GB mensuales de \textit{logs} y otros tantos de trazas con catorce días de retención,\footnote{Precios oficiales: \url{https://grafana.com/pricing/}} \textbf{New Relic} concede 100 GB mensuales de datos de cualquier origen y un usuario con acceso completo,\footnote{Precios oficiales: \url{https://newrelic.com/pricing/free-tier}} y \textbf{Axiom} llega a 500 GB mensuales con treinta días de retención.\footnote{Precios oficiales: \url{https://axiom.co/pricing}} Una tercera vía consistiría en desplegar una solución propia, como \textbf{Grafana} junto a \textbf{Loki} y \textbf{Tempo}, sobre infraestructura del proyecto.

Esta última se descartó por la misma razón que condujo a elegir Lambda para la API: exige un proceso en ejecución permanente. El servidor más reducido capaz de sostener un sistema de este tipo consumiría por sí solo la totalidad del margen de gasto fijado para el proyecto, de modo que observar el sistema costaría más que ejecutarlo. Sostener una infraestructura que no factura mientras no se usa mediante otra que factura por horas anularía precisamente la ventaja que motivó la primera decisión.

Las plataformas externas se descartaron por dos motivos que se refuerzan entre sí. El primero es técnico: una función de Lambda no mantiene un proceso vivo entre invocaciones, por lo que enviar la telemetría a un destino externo obliga a hacerlo dentro de la propia petición —lo que añade latencia a cada respuesta— o a desplegar junto a la función un componente adicional que la recoja y la reenvíe. X-Ray no plantea ese problema, porque el entorno de ejecución de Lambda incluye un \textit{daemon} que recibe la telemetría de forma local y se encarga de enviarla.\footnote{Variables de entorno de Lambda: \url{https://docs.aws.amazon.com/lambda/latest/dg/configuration-envvars.html}} El segundo es económico: Lambda escribe sus \textit{logs} en CloudWatch sin que el destino sea configurable, de manera que ese volumen se consume en cualquier caso; remitirlo además a una plataforma externa supondría mantener dos sistemas de observabilidad para un único conjunto de datos.

Se optó en consecuencia por \textbf{CloudWatch} y \textbf{X-Ray}. Su capa gratuita cubre con holgura el tráfico previsto: 5 GB mensuales de datos de \textit{logs} —sumando la ingesta, el almacenamiento y el volumen examinado por las consultas—, 100.000 trazas registradas al mes, diez alarmas y tres paneles de hasta cincuenta métricas cada uno; superados esos límites, la ingesta se factura a 0,50 USD por GB y el almacenamiento a 0,03 USD por GB.\footnote{Precios oficiales: \url{https://aws.amazon.com/cloudwatch/pricing/}} Igual que se argumentó al descartar API Gateway, lo determinante no es la amplitud de la capa gratuita sino su permanencia: la de CloudWatch y X-Ray no expira a los doce meses de creada la cuenta, por lo que adoptarlas no programa un gasto futuro.

Elegido el destino, quedaba decidir cómo genera la API las trazas que lo alimentan. El rastreo activo de Lambda registra por sí solo cada invocación y separa la inicialización del entorno del tiempo de ejecución, pero no reparte este último entre las operaciones de la aplicación, que es lo que pide \reqref{RNF-MO-03}. La vía tradicional para obtener ese desglose, el SDK de X-Ray, entró en modo de mantenimiento en febrero de 2026 y dejará de recibir soporte en febrero de 2027, y \textbf{AWS} recomienda sustituirlo por \textbf{OpenTelemetry}, cuyas trazas X-Ray recibe sin cambios en la consola.\footnote{Calendario de fin de soporte del SDK de X-Ray: \url{https://docs.aws.amazon.com/xray/latest/devguide/xray-daemon-eos.html}} Adoptar en un proyecto nuevo una dependencia con fecha de retirada habría obligado a migrarla pocos meses después.

Se optó por el \textit{layer} de OpenTelemetry que publica \textbf{AWS} para Lambda, que aporta el SDK y lo configura para enviar las trazas al \textit{daemon} de X-Ray del entorno de ejecución.\footnote{Instrumentación de Lambda con OpenTelemetry: \url{https://docs.aws.amazon.com/xray/latest/devguide/migrate-xray-to-opentelemetry-python.html}} La aplicación no contiene instrumentación propia: \textbf{FastAPI} emite de forma nativa, desde su versión 0.142, un \textit{span} por petición con el nombre de la ruta que la atiende y otro por cada fase de su procesamiento —la resolución de dependencias, la ejecución de la operación y la serialización de la respuesta—, y el \textit{layer} los incorpora a la traza de la invocación. El precio de esta solución es el arranque en frío, que se alarga porque el SDK se importa antes que la aplicación.

\begin{figure}[htbp]
	\centering
	\includegraphics[width=0.85\textwidth]{monitoring-infrastructure.pdf}
	\caption{Infraestructura de la monitorización.}
	\label{fig:monitoring-infrastructure}
\end{figure}

El diseño satisface \reqref{RNF-MO-01} mediante el registro de los errores, \reqref{RNF-MO-03} mediante las trazas, que desglosan el tiempo de cada petición por operación, y de forma parcial \reqref{RNF-MO-04}, ya que el recuento de invocaciones mide peticiones atendidas y no usuarios activos. Queda sin satisfacer \reqref{RNF-MO-02}, por un motivo económico: la capa gratuita que cubriría las notificaciones es acotada y nada limita por sí solo el número de avisos que un sistema de alertas emite, de modo que permanecer dentro de ella exigiría añadir mecanismos de contención que son a su vez infraestructura capaz de fallar. Al tratarse de una prioridad \textit{Should Have}, la decisión es reversible y se revisará cuando el sistema reciba tráfico real.

El modelo presenta, por último, dos puntos ciegos que conviene hacer explícitos. El primero es CloudFront: la distribución no participa en las trazas, de modo que el tiempo transcurrido entre la petición del navegador y la invocación de la función no queda medido y la traza comienza en la función misma. El segundo es el propio navegador, ya que el comportamiento del front-end en el cliente —tiempos de carga y errores de JavaScript— no se recoge: el servicio de AWS destinado a ello, \textbf{CloudWatch RUM}, se factura a 1 USD por cada 100.000 eventos y su gratuidad se limita a una prueba inicial de un millón de eventos por cuenta, no a una asignación mensual recurrente.
